\section{Conclusion}

We presented a context-aware authoring workflow that exposes an AI-proposed shared dyadic interpretation through editable Semantic Performance Tags between narrative interpretation and procedural character animation. A pretrained LLM proposes initial performance states for both interlocutors and sparse dialogue-anchored semantic changes, while the animator can inspect and revise one or multiple executable decisions across either character, including affect, attention, head behavior, and performative blink before deterministic compilation of the current canonical plan to JALI and Maya. The representation separates semantic acting decisions from shot- and rig-dependent execution, including calibrated realization of semantic gaze targets.

In the context ablation, Full-Context proposals were rated significantly higher on overall scene fit than Dialogue-Only proposals, while individual intention, affect, gaze, and timing ratings did not show corrected significant differences. In the first-use authoring comparison, editable Semantic Performance Tags produced directional improvements in final-performance ratings and revision clarity but no corrected significant workflow or workload advantage. The current evidence therefore supports a cautious interpretation: contextual information contributes to holistic performance coherence, and editable semantic tags can function as an inspectable authoring surface, but this N=10 snapshot does not establish a general superiority over prompt-level regeneration.

The broader question is not whether an AI system can generate a plausible facial performance in one step, but how a creator can inspect and negotiate the interpretation that produced it. Our results motivate treating performance interpretation as an editable layer between narrative context and deterministic execution, while leaving claims about efficiency, long-term authorship, and detailed accept/modify/reject behavior to larger and more fully instrumented studies.
